% Standalone fragment: requires the host's suppthm environment and amsmath.
% Adapted from the September 18 contextual companion
% companions/mistake_deadline_frontier.tex. The Pareto-minimality argument
% below strengthens its schedule-domination statement.
\subsection{Optimal mistake and convergence guarantees for the staircase family}
\label{app:p29-frontier}

In this subsection we prove Theorem~\ref{thm:staircase-frontier}, which
determines the optimal mistake and convergence guarantees for the staircase
family. This family generalizes the construction in Theorem~6.4
of~\citet[][P29]{mistake} by allowing arbitrary positive finite block sizes.
We compare generators on every possible target simultaneously: improving
one target's guarantee may worsen another's. The proof constructs generators
from a sequence of binary choices, derives matching lower bounds for
arbitrary generators, and shows that every resulting choice expresses an
optimal tradeoff.

\subsubsection{Languages, histories, and guarantees}
\label{app:p29-frontier-setup}
We first give a concrete realization of the family in
Section~\ref{sec:new-mathematics}. Fix positive integers $a_1,a_2,\ldots$. On
$X=\mathbb N_{>0}\times\mathbb N_{>0}$, put
\[
 C_i=\{(i,j):1\le j\le a_i\},\qquad
 P_i=\bigcup_{k=1}^{i}C_k,\qquad
 s_i=|P_i|=\sum_{k=1}^{i}a_k,
 \qquad s_0=0.
\]
Write $R_i=\{(i,j):j>a_i\}$ for the private remainder of column $i$,
and define
\[
 L_i=P_i\cup R_i
     =P_{i-1}\cup(\{i\}\times\mathbb N_{>0}).
\]
Each $L_i$ is infinite. For $k<i$, its intersection with $L_k$ is
$P_k$. A point of $C_k$ belongs to exactly the languages $L_i$ with
$i\ge k$, whereas a point of $R_k$ belongs only to $L_k$.
For the abstract family in the main text, choose a bijection from each
finite block to the first $a_i$ points of its column and from its countably
infinite private part to the remainder. These bijections preserve every
language membership. Any additional ambient points lie in no target;
the lower-bound argument below also allows a generator to output them.
Thus the concrete notation entails no restriction on the stated family.

A legal history for target $L_i$ is a finite ordered list of distinct
points in $L_i$. After a history of length $t\ge0$, a generator $G$
outputs a point outside its observed set, before the next observation
arrives. A generator is a deterministic function of the entire ordered
history; no computability restriction is imposed. An output is
a mistake when it lies outside the fixed target. This definition imposes
freshness relative to observed inputs, without requiring the outputs to
be mutually distinct.

Let $M_i(G)\in\mathbb N\cup\{\infty\}$ be the supremum of the total
number of mistakes over injective positive streams in $L_i$. Let
$D_i(G)\in\mathbb N\cup\{\infty\}$ be the least nonnegative integer
$D$ such that $G$ is valid on every legal $L_i$-history of every length
$t\ge D$; set $D_i(G)=\infty$ if no such finite $D$ exists. In
particular, an error after exactly $s$ observations forces
$D_i(G)\ge s+1$, and $D_i(G)=0$ means that no legal history receives
an invalid output. Time is the number of distinct observations,
including time zero at the empty history.

The same worst-case quantities result if streams are required to
enumerate the entire target. Every finite injective positive history
extends to such an enumeration. Thus every finite collection of
mistakes witnessed on a positive stream is also witnessed on some
complete text, and every history used to test a convergence time occurs on a
complete text.

The two quantities answer different questions: $M_i(G)$ bounds how many
errors can occur, while $D_i(G)$ bounds how late an error can occur.
We compare both quantities for all targets at once. Specifically, $G'$ dominates $G$ if
$M_i(G')\le M_i(G)$ and $D_i(G')\le D_i(G)$ for every $i$. A profile
is Pareto minimal if no generator dominates it with a strict inequality
in at least one coordinate. Thus a Pareto-minimal profile cannot be improved
on any target without worsening at least one of the guarantees.

\subsubsection{Characterization and proof}
\label{app:p29-frontier-proof}
The relevant decisions occur when the observed set is exactly $P_k$.
At that point the possible targets are $L_k,L_{k+1},\ldots$, whose
intersection is $P_k$, so no unobserved point is valid for all of them.
Choosing a point in $R_k$ is correct only for $L_k$; choosing a point
in $C_{k+1}$ is correct for every later target. We record these choices
by a binary sequence $\sigma=(\sigma_k)_{k\ge1}$, with $\sigma_k=1$
for the first choice and $\sigma_k=0$ for the second. The construction
below makes no mistakes away from these decision points.

For target $L_i$, an earlier choice of $R_k$, with $k<i$, causes one
mistake if the observed set reaches $P_k$. The choice of $C_{i+1}$
causes a mistake at $P_i$. These observations motivate the candidate
mistake bound and the corresponding set of possible error times:
\[
 m_i(\sigma)=\sum_{k<i}\sigma_k+1-\sigma_i,
 \qquad
 E_i(\sigma)=\{s_k:k<i,\ \sigma_k=1\}
              \cup\{s_i:\sigma_i=0\}.
\]
The proposed convergence time is one more than the latest of these times, or
zero if there are none:
\begin{equation}
\label{app:p29-frontier-deadline}
 d_i(\sigma)=
 \begin{cases}
  0,&E_i(\sigma)=\varnothing,\\
  1+\max E_i(\sigma),&E_i(\sigma)\ne\varnothing.
 \end{cases}
\end{equation}

The first part of the theorem shows that these formulas are attained.
The second proves that every generator has guarantees at least as large
as those of one such sequence, simultaneously for every target. The third
establishes that none of the resulting profiles can be improved further.

\begin{suppthm}[Optimal mistake and convergence guarantees]
\label{app:p29-frontier-theorem}
For the staircase family above, the following assertions hold.
\begin{enumerate}
\item For every binary sequence $\sigma$, there is a deterministic
generator $G_\sigma$ satisfying, simultaneously for all $i\ge1$,
\begin{equation}
\label{app:p29-frontier-exact}
 M_i(G_\sigma)=m_i(\sigma),\qquad
 D_i(G_\sigma)=d_i(\sigma).
\end{equation}
\item Every deterministic generator $G$ is dominated by a schedule:
there is one $\sigma$ such that, for all $i\ge1$,
\begin{equation}
\label{app:p29-frontier-domination}
 m_i(\sigma)\le M_i(G),\qquad d_i(\sigma)\le D_i(G).
\end{equation}
\item Every schedule profile is Pareto minimal. Consequently the
profiles in~\eqref{app:p29-frontier-exact} are exactly the deterministic
Pareto frontier.
\end{enumerate}
\end{suppthm}
\begin{proof}
\emph{Attaining the proposed guarantees.}
We first define $G_\sigma$ on every finite history and then prove the
two equalities in~\eqref{app:p29-frontier-exact}.
Choose fixed orders for making all the selections below. If a history
contains a point of $R_i$, any compatible target must be $L_i$.
At such a realizable history, output any unobserved point of $L_i$.
If a history is unrealizable, output any unobserved point of $X$;
this branch is irrelevant to targetwise guarantees but makes the
generator total and fresh.

It remains to define the generator before a private point appears.
At the empty history, output a point of $C_1$, which belongs to every
target. At a nonempty history consisting only of finite-block points,
let $k$ be the largest observed column index and let $S$ be its
observed set. The compatible targets are precisely
$L_k,L_{k+1},\ldots$, and their intersection is $P_k$. Indeed, the
observed point of $C_k$ excludes every earlier target, while every
target of index at least $k$ contains all the observed points. Their
common core is $P_k$ because $L_k\cap L_{k+1}=P_k$.

If $S\ne P_k$, output a point of $P_k\setminus S$. It is fresh and
valid for every compatible target. If $S=P_k$, no fresh point is valid
for all compatible targets. When $\sigma_k=1$, output a point of
$R_k$, valid only for $L_k$. When $\sigma_k=0$, output a point of
$C_{k+1}$, valid for every later target and invalid for $L_k$. Both
choices are fresh.

Fix a target $L_i$. Its stream can encounter a boundary $S=P_k$
only for $k\le i$, and it can encounter each such boundary at most
once because observations are injective. Every output away from
these boundaries is valid. At a boundary $k<i$, a mistake occurs
exactly when $\sigma_k=1$; at boundary $k=i$, a mistake occurs
exactly when $\sigma_i=0$. The respective observed counts are $s_k$.
This proves $M_i(G_\sigma)\le m_i(\sigma)$ and
$D_i(G_\sigma)\le d_i(\sigma)$.

To attain these bounds, present $C_1,C_2,\ldots,C_i$ in fixed
internal orders, followed by an enumeration of $R_i$. This is an
injective complete text of $L_i$, and it encounters every boundary
$P_k$ for $k\le i$. It therefore realizes exactly the mistakes in
$E_i(\sigma)$. If that set is nonempty, its latest mistake occurs
after $\max E_i(\sigma)$ observations, forcing convergence time at least
$1+\max E_i(\sigma)$. If it is empty, the preceding upper bound
already gives convergence time zero. This establishes both equalities
in~\eqref{app:p29-frontier-exact}.

\emph{Lower bounds for an arbitrary generator.}
We next show that allowing more elaborate choices, including dependence
on the order of observations, cannot improve on all the profiles above.
Fix a deterministic fresh generator $G$. For each $k$, let $h_k$
be the ordered concatenation of $C_1,\ldots,C_k$ using the same
fixed internal orders. Define
\[
 \sigma_k=1\quad\Longleftrightarrow\quad G(h_k)\in R_k.
\]
Freshness gives $G(h_k)\notin P_k$. Hence if $\sigma_k=0$,
the output at $h_k$ is outside $L_k=P_k\cup R_k$.

For a fixed target $L_i$, extend $h_i$ by the private remainder
$R_i$ to obtain one complete $L_i$-text. At every earlier boundary
$h_k$ with $\sigma_k=1$, the output belongs to $R_k$ and is
invalid for $L_i$. At its own boundary $h_i$, the output is invalid
if $\sigma_i=0$. Thus this one text witnesses all
$m_i(\sigma)$ charged mistakes, at exactly the observed counts
in $E_i(\sigma)$. The resulting mistake and convergence time lower bounds
are~\eqref{app:p29-frontier-domination}; other errors made by $G$
can only increase those quantities.

The infinite concatenation of all finite blocks is used only to
specify a consistent family of finite histories. It is not asserted
to be a text of any one target. Each lower-bound argument above
fixes $i$ and uses the legal prefix $h_i$ followed by a legal
continuation for that same target.

\emph{Optimality of every schedule.}
The lower bound shows that the schedules suffice to find all optimal
guarantees. We must still prove that every schedule is itself optimal.
First, no two distinct schedules have coordinatewise ordered
mistake vectors. Suppose otherwise that
$m_i(\tau)\le m_i(\sigma)$ for every $i$, and let $k$ be the
first index where $\tau_k\ne\sigma_k$. The inequality at $i=k$
forces $\tau_k=1$ and $\sigma_k=0$. At $i=k+1$, the difference
of the mistake counts is
\[
 m_{k+1}(\tau)-m_{k+1}(\sigma)
 =1-(\tau_{k+1}-\sigma_{k+1}).
\]
For this to be nonpositive, we must also have
$\tau_{k+1}=1$ and $\sigma_{k+1}=0$. But then
\[
 m_{k+2}(\tau)-m_{k+2}(\sigma)
 =2-(\tau_{k+2}-\sigma_{k+2})\ge1,
\]
a contradiction. Therefore $\tau=\sigma$ whenever all the
mistake inequalities hold.

Now suppose an arbitrary $G$ dominates $G_\sigma$. The reduction
just proved supplies a schedule $\tau$ with
$m_i(\tau)\le M_i(G)\le m_i(\sigma)$ for every $i$, so
$\tau=\sigma$. Applying the convergence time part of the same reduction
gives $d_i(\sigma)\le D_i(G)\le d_i(\sigma)$, while the mistake
inequalities also become equalities. No coordinate can improve
strictly. Thus every schedule profile is Pareto minimal. Conversely,
every Pareto-minimal generator must have the same profile as its
dominating schedule, proving the final assertion.
\end{proof}

The two constant schedules illustrate the tradeoff. Always choosing
$C_{k+1}$ makes no mistakes at the earlier decision points for target
$L_i$, but makes one at $P_i$. Its guarantees are therefore $M_i=1$
and $D_i=s_i+1$. Always choosing $R_k$ avoids the mistake at $P_i$
but makes one at each earlier decision point, giving $M_i=i-1$ and
$D_i=s_{i-1}+1$ for $i\ge2$, with $M_1=D_1=0$.
Theorem~\ref{app:p29-frontier-theorem} determines all intermediate
choices as well. Its lower bound applies to every deterministic generator,
and its minimality statement shows that each listed profile is an optimal
tradeoff across the targets, completing the special-case characterization.
